% !TeX root = ../main.tex
\chapter{KẾT LUẬN}\label{chap:conclusion}

\section{Tổng kết}
Khóa luận đã giải quyết bài toán xây dựng nền tảng đa doanh nghiệp hỗ trợ tư vấn khách hàng đa kênh, quản lý đơn hàng và điều phối quy trình bàn giao nội bộ có kiểm soát trí tuệ nhân tạo cho các doanh nghiệp vừa và nhỏ kinh doanh sản phẩm vật lý. Xuất phát từ những hạn chế thực tiễn tại Doanh nghiệp A và khoảng trống thiết kế giữa ba nhóm giải pháp hiện có trên thị trường (nền tảng bán lẻ -- kho, công cụ gom hội thoại đa kênh và hệ thống ERP tổng thể), khóa luận đạt được bốn nhóm đóng góp chính từ phân tích, thiết kế đến cài đặt và kiểm chứng:

\begin{enumerate}
  \item \textbf{Phân tích bài toán và phân định ranh giới phạm vi minh bạch:} Hệ thống hóa bốn thách thức cốt lõi từ khâu tiếp thị đầu phễu đến tư vấn, đơn nháp và bàn giao trách nhiệm; xác lập bộ yêu cầu chức năng và phi chức năng chi tiết, đồng thời phân tách rõ ràng giữa phần hạt nhân đã cài đặt/kiểm chứng trong môi trường cô lập và các hạng mục mở rộng ở mức đặc tả mục tiêu (\emph{TARGET}) nhằm bảo đảm tính trung thực học thuật.
  \item \textbf{Thiết kế kiến trúc đa dịch vụ phân tách trách nhiệm:} Xây dựng kiến trúc phân tách rõ thẩm quyền giữa bốn thành phần chính: lớp giao diện người dùng Next.js (\texttt{apps/\hspace{0pt}web} và \texttt{packages/\hspace{0pt}sdk}); dịch vụ nghiệp vụ trung tâm bằng Go (\texttt{services/\hspace{0pt}api} cùng tiến trình tiêu thụ sự kiện \texttt{cmd/\hspace{0pt}ingress}) nắm độc quyền xác thực phiên trên cơ sở dữ liệu, kiểm chứng quyền theo trạng thái hiện hành và phát lệnh nghiệp vụ; dịch vụ kết nối kênh trực tiếp bằng Go (\texttt{services/\hspace{0pt}connector} với hai bộ điều hợp \texttt{plugin/\hspace{0pt}facebook} và \texttt{plugin/\hspace{0pt}telegram}) cô lập giao tiếp webhook/\hspace{0pt}API bên ngoài, mã hóa khóa kênh bằng AES-256-GCM gắn ngữ cảnh AAD và trao đổi bất đồng bộ qua Apache Kafka (Redpanda) theo mẫu \emph{Transactional Outbox / Consumer Inbox}; và dịch vụ thực thi AI bằng Python (\texttt{services/\hspace{0pt}ai} qua gRPC/\hspace{0pt}Protobuf với LangGraph và bộ điều hợp LangChain đa nhà cung cấp) chỉ tiếp nhận ảnh chụp ngữ cảnh được cấp phép để trả về ứng viên, tuyệt đối không có quyền tự ghi dữ liệu nghiệp vụ hay phát lệnh gửi tin.
  \item \textbf{Thiết kế cơ sở dữ liệu 39 bảng trên ba vùng sở hữu Atlas độc lập:} Tổ chức dữ liệu bền vững trên PostgreSQL~17 thành ba vùng sở hữu tách biệt được quản lý tiến hóa tường minh bằng 22 tệp migration Atlas (chi tiết tại Phụ lục~\ref{app:database}): lược đồ \texttt{public} gồm 31 bảng thuộc \texttt{hexta\_\hspace{0pt}api}, lược đồ \texttt{ai\_\hspace{0pt}runtime} gồm 1 bảng (\texttt{ai\_\hspace{0pt}runtime.\hspace{0pt}runs}) thuộc \texttt{hexta\_\hspace{0pt}ai} và lược đồ \texttt{connector} gồm 7 bảng thuộc \texttt{hexta\_\hspace{0pt}connector}. Thiết kế áp dụng khóa ngoại tổ hợp \texttt{(tenant\_id, id)} ở các cụm quan hệ trọng yếu, kiểm soát cạnh tranh lạc quan theo \texttt{version}, cờ \texttt{price\_\hspace{0pt}provided} phân biệt chưa nhập giá với giá bằng không, chỉ mục duy nhất bộ phận (\texttt{idx\_\hspace{0pt}messaging\_\hspace{0pt}one\_\hspace{0pt}open\_\hspace{0pt}turn}, \texttt{uq\_\hspace{0pt}messaging\_\hspace{0pt}connections\_\hspace{0pt}active\_\hspace{0pt}page}) và cơ chế thuê khóa \texttt{lease} chống gọi lặp mô hình có phí.
  \item \textbf{Cài đặt và kiểm chứng trên môi trường thực nghiệm cô lập:} Hoàn thiện cài đặt lõi cùng bộ kiểm thử tự động trong môi trường Docker Compose có PostgreSQL~17 và Redpanda cô lập: bộ kiểm thử Go với \texttt{-race} vượt qua toàn bộ 19 gói thuộc \texttt{services/\hspace{0pt}api}, \texttt{services/\hspace{0pt}connector} và \texttt{packages/\hspace{0pt}shared} (bao gồm các ca đầu-cuối cho Facebook OAuth, Webhook, Send, Telegram Bot và \texttt{Outbox\hspace{0pt}Relay}); dịch vụ AI Python vượt qua 109 bài kiểm thử (0 bài bỏ qua); gói \texttt{@hexta/\hspace{0pt}sdk} vượt qua 35 bài kiểm thử; các kiểm tra tĩnh và đóng gói (\texttt{gofmt}, \texttt{go vet}, \texttt{make migrate-\hspace{0pt}validate}, \texttt{make swagger}, kiểm tra lint/build giao diện web bằng \texttt{pnpm}, triển khai thử cấu hình \texttt{compose.\hspace{0pt}prod.\hspace{0pt}yml} trên \texttt{localhost} và 9 kịch bản tác nhân \texttt{scripts/\hspace{0pt}deploy-\hspace{0pt}agent.\hspace{0pt}sh} trong container Debian) đều đạt kết quả thành công.
\end{enumerate}

Đối chiếu với ba câu hỏi nghiên cứu đặt ra ở Chương~\ref{chap:introduction}, khóa luận thu được các kết luận tổng hợp sau:

\begin{itemize}
  \item \textbf{Đối với RQ1 (Điều phối quy trình tuần tự và trách nhiệm bàn giao):} Mô hình tách \texttt{workflow\_\hspace{0pt}definitions}, \texttt{workflow\_\hspace{0pt}versions} và \texttt{workflow\_\hspace{0pt}steps} (tối đa 30 bước) kết hợp hai chế độ bàn giao \texttt{direct} và \texttt{acceptance} cho phép mỗi doanh nghiệp tùy biến chuỗi xử lý phù hợp với cơ cấu nhân sự của mình mà không làm xáo trộn công việc đang vận hành khi cấu hình thay đổi. Tại mọi thời điểm, trường \texttt{current\_\hspace{0pt}step}, \texttt{responsible\_\hspace{0pt}id}, \texttt{recipient\_\hspace{0pt}id} và trạng thái (\texttt{active}, \texttt{waiting}, \texttt{completed}) xác định rõ bước hiện hành và người chịu trách nhiệm; tranh chấp giữa lệnh \texttt{accept} và \texttt{assign} được tuần tự hóa an toàn ở tầng giao dịch và lưu vết đầy đủ trong \texttt{work\_\hspace{0pt}events}.
  \item \textbf{Đối với RQ2 (Định tuyến kịch bản/AI có kiểm chứng và hàng rào điều khiển Desk):} Đồ thị LangGraph kiểm soát chặt đầu ra có cấu trúc, đối chiếu trích dẫn theo dữ liệu doanh nghiệp được cấp phép và tự động chuyển nhân viên khi thiếu thông tin; trong luồng tư vấn mua sắm SSE (\texttt{/api/\hspace{0pt}v1/\hspace{0pt}ai/\hspace{0pt}converse}), máy chủ Go đối chiếu lại danh mục tối đa 100 sản phẩm đang hoạt động có giá và tự tính lại thành tiền để loại bỏ nguy cơ AI bịa đặt giá bán. Đồng thời, phân hệ Desk quản lý ba trạng thái chiếu (\texttt{BOT\_\hspace{0pt}ACTIVE}, \texttt{WAITING\_\hspace{0pt}STAFF}, \texttt{STAFF\_\hspace{0pt}ACTIVE}) đi kèm bộ đếm đơn điệu \texttt{control\_\hspace{0pt}generation} và danh mục lý do chuyển giao chuẩn hóa, bảo đảm chặn được thao tác dựa trên ngữ cảnh cũ khi nhân viên tiếp nhận. Về mặt đo đạc định lượng thời gian thao tác thực tế của nhân viên với mô hình ngôn ngữ lớn thật, phạm vi đánh giá trong khóa luận mới kiểm chứng mức hợp đồng, cấu trúc và rào chắn; thực nghiệm bấm giờ trên tập dữ liệu đầy đủ được ghi nhận rõ là một hạn chế mở cần tiếp tục thực hiện.
  \item \textbf{Đối với RQ3 (Cô lập đa doanh nghiệp, kiểm soát cạnh tranh và lưu vết):} Hệ thống bảo đảm ranh giới đa doanh nghiệp thông qua cơ chế xét duyệt đăng ký và lời mời băm SHA-256, kiểm chứng tư cách thành viên hiện hành trực tiếp trên cơ sở dữ liệu (\texttt{Require\hspace{0pt}Tenant\hspace{0pt}Member}, \texttt{Authorize\hspace{0pt}Member}), khóa ngoại tổ hợp kèm \texttt{tenant\_\hspace{0pt}id}, khóa hàng kết hợp kiểm tra \texttt{version}, \texttt{context\_\hspace{0pt}revision} và \texttt{control\_\hspace{0pt}generation} chống cập nhật dựa trên dữ liệu cũ, khóa chống trùng lặp (\texttt{Idempotency-\hspace{0pt}Key}, \texttt{request\_\hspace{0pt}key}, \texttt{client\_\hspace{0pt}message\_\hspace{0pt}id}) và chuỗi lưu vết đồng giao dịch tại \texttt{audit\_\hspace{0pt}logs}, \texttt{work\_\hspace{0pt}events}, \texttt{domain\_\hspace{0pt}events} và \texttt{messaging\_\hspace{0pt}ingress}.
\end{itemize}

\section{Hạn chế của khóa luận}
Bên cạnh các kết quả đã cài đặt và kiểm chứng cô lập, khóa luận nhìn nhận thẳng thắn bốn hạn chế sau:

\begin{enumerate}
  \item \textbf{Chưa có thực nghiệm bấm giờ nhân viên cho RQ2 với khóa mô hình thật và bộ dữ liệu 10 cây hoàn chỉnh:} Việc đánh giá phân hệ AI hiện tại chủ yếu dựa trên các ca kiểm thử hợp đồng gRPC/\hspace{0pt}HTTP, kiểm chứng cơ chế thuê khóa \texttt{lease}, hàng rào \texttt{control\_\hspace{0pt}generation}, kiểm tra định dạng đầu ra có cấu trúc và bộ điều hợp mô phỏng (\texttt{fake}) hoặc giả lập truyền tải LangChain. Khóa luận chưa tiến hành thực nghiệm bấm giờ đối chứng thời gian chuẩn bị, rà soát và chỉnh sửa phản hồi của nhân viên tư vấn trên bộ dữ liệu 10 loại cây trồng hoàn chỉnh với khóa API mô hình ngôn ngữ lớn thật.
  \item \textbf{Kết nối Facebook Page và Telegram Bot chưa nghiệm thu trên Internet công cộng:} Các kịch bản OAuth, đăng ký webhook, tiếp nhận tin nhắn, xử lý phần giữ chỗ phương tiện không hỗ trợ và gửi tin một lần của \texttt{services/\hspace{0pt}connector} mới được nghiệm thu đầy đủ trước máy chủ giả lập giao thức Meta Graph API và Telegram Bot API cục bộ, kèm kiểm thử gọi lẻ API \texttt{getMe} và \texttt{sendMessage} lỗi (\texttt{chat\_\hspace{0pt}not\_\hspace{0pt}found}, \texttt{auth\_\hspace{0pt}expired}) tới Telegram. Việc nghiệm thu webhook HTTPS công khai qua Internet gắn với quy trình \emph{Meta App Review} và kiểm thử tương tác giao diện trình duyệt toàn trình hai chiều vẫn chưa được thực thi (\emph{NOT\_RUN}).
  \item \textbf{Một số hạng mục mở rộng trong thiết kế tổng thể mới ở dạng đặc tả mục tiêu:} Để tập trung hoàn thiện trục cốt lõi trong khuôn khổ thời gian khóa luận, các thành phần đã phân tích trong Chương~\ref{chap:design} nhưng chưa được hiện thực hóa thành mã chạy gồm: phân hệ kịch bản có phiên bản (\texttt{scripts}/\hspace{0pt}\texttt{script\_\hspace{0pt}versions}) với bốn quyền độc lập do chủ doanh nghiệp cấp (\texttt{script.\hspace{0pt}create}, \texttt{script.\hspace{0pt}edit}, \texttt{script.\hspace{0pt}activate}, \texttt{script.\hspace{0pt}archive}); bộ lập lịch tự động đóng lượt tư vấn sau 24 giờ không hoạt động (ngoại trừ trạng thái \texttt{needs\_\hspace{0pt}human}); quy trình xử lý ngoại lệ hủy đơn và trả hàng (\texttt{exception\_\hspace{0pt}cases}/\hspace{0pt}\texttt{return\_\hspace{0pt}lines}); cùng bộ quét tự động thu hồi các bản ghi outbox vượt quá 12 lần thử (khi cạn số lần thử, ý định gửi hiện vẫn lưu trạng thái \texttt{pending}/\hspace{0pt}\texttt{sending}) và cơ chế kiểm tra định kỳ tình trạng hiệu lực token kênh ngoài thời điểm gửi tin.
  \item \textbf{Chưa thực hiện kiểm thử tải quy mô lớn và đánh giá cụm đa nút:} Toàn bộ quá trình thực nghiệm mới diễn ra trên môi trường cục bộ và cấu hình triển khai máy chủ đơn (\texttt{infrastructure/\hspace{0pt}compose.\hspace{0pt}prod.\hspace{0pt}yml} thiết kế cho máy chủ 2~vCPU, 16~GB RAM dùng Docker Compose và Caddy), chưa thực hiện kiểm thử tải đồng thời cao để đo thông lượng giới hạn của Kafka/\hspace{0pt}Redpanda, độ trễ truy vấn PostgreSQL khi dữ liệu tích lũy lớn, diễn tập sao lưu -- phục hồi thực địa hay đánh giá hiệu năng trên cụm đa nút (\emph{multi-node}).
\end{enumerate}

\section{Định hướng phát triển}
Từ nền tảng kiến trúc và ranh giới đã xác lập, hệ thống có ba hướng phát triển tiếp theo sau khóa luận:

\begin{enumerate}
  \item \textbf{Nghiệm thu thực địa hạ tầng kênh thật và thực nghiệm đối chứng thao tác nhân viên:} Triển khai cấu hình \texttt{infrastructure/\hspace{0pt}compose.\hspace{0pt}prod.\hspace{0pt}yml} cùng tác nhân cập nhật \texttt{scripts/\hspace{0pt}deploy-\hspace{0pt}agent.\hspace{0pt}sh} lên máy chủ thực tế có tên miền và chứng chỉ TLS công cộng; hoàn tất kiểm thử webhook thực địa cho Telegram Bot và Facebook Page sau khi có quyền Meta App Review; đồng thời thực hiện thực nghiệm đối chứng đo thời gian thao tác và tỷ lệ chỉnh sửa phản hồi của nhân viên trên bộ dữ liệu 10 cây hoàn chỉnh với khóa mô hình thật, kết hợp kiểm thử tải và diễn tập sao lưu -- khôi phục dữ liệu.
  \item \textbf{Hoàn thiện các cơ chế điều phối tự động, đối soát và quan sát AI:} Bổ sung bộ lập lịch đóng lượt tư vấn theo quy tắc 24 giờ không hoạt động tôn trọng cờ \texttt{needs\_\hspace{0pt}human}; hoàn thiện luồng đối soát tự động trạng thái gửi \texttt{unknown}, bộ quét xử lý bản ghi outbox quá 12 lần thử và tác vụ kiểm tra định kỳ hiệu lực token kết nối; triển khai đầy đủ cụm bảng \texttt{scripts}/\hspace{0pt}\texttt{script\_\hspace{0pt}versions} với bốn quyền độc lập và quy trình ngoại lệ hủy/\hspace{0pt}trả hàng; đồng thời tích hợp công cụ quan sát luồng thực thi AI tự chủ (như Langfuse tự lưu trữ) phục vụ theo dõi chi phí và chất lượng gợi ý.
  \item \textbf{Mở rộng hệ sinh thái đa ứng dụng Hexta -- PostPilot -- ViralForge:} Theo lộ trình phát triển hệ sinh thái sau cột mốc khóa luận, kiến trúc mở với biên định danh doanh nghiệp dùng chung (\texttt{tenants}, \texttt{tenant\_\hspace{0pt}members}) và trục sự kiện Kafka sẽ kết nối ba ứng dụng có vai trò phân tách rõ ràng: \textbf{Hexta} giữ vai trò trung tâm vận hành đơn hàng, tư vấn khách hàng và điều phối công việc nội bộ; \textbf{PostPilot} đảm nhiệm quản trị kết nối trang mạng xã hội, lập lịch và xuất bản nội dung tiếp thị; \textbf{ViralForge} đảm nhiệm khởi tạo nội dung truyền thông đa phương tiện có kiểm soát chi phí và quyền sở hữu. Hai ứng dụng vệ tinh sẽ được hoàn thiện lớp cô lập đa doanh nghiệp trước khi liên thông luồng bàn giao tài nguyên từ ViralForge sang PostPilot và đồng bộ dữ liệu báo cáo hợp nhất về Hexta.
\end{enumerate}
